\subsection{Risoluzione dei problemi}
\label{subsection:risoluzione_problemi}

\subsubsection{Scopo}
La risoluzione dei problemi è un processo che ha lo scopo di individuare e risolvere le problematiche che si presentano durante lo svolgimento del progetto.
Il gruppo deve controllare costantemente lo stato di avanzamento delle attività e, in caso di problemi, deve attivarsi per risolverli nel minor tempo possibile.
Ciò deve essere fatto andando a documentarne le cause e le soluzioni adottate, in modo da evitare che si ripresentino in futuro.

\subsubsection{Esecuzione}
Quando viene riscontrata una problematica all'interno dei prodotti, i membri del gruppo devono segnalarla tempestivamente all'interno del \glossario{repository} del progetto con l'etichetta "bug", in modo che possa essere risolta nel minor tempo possibile.
Deve essere inoltre notificato il gruppo tramite i \hyperref[subpar:canali_interni]{canali di comunicazione interni}.
Il o i membri del gruppo che rilevano il problema devono descriverlo in modo chiaro e dettagliato, in modo che chi dovrà risolverlo possa capire immediatamente di cosa si tratta.
Se invece la problematica venisse riscontrata durante una riunione interna, questa dovrà essere documentata all'interno del verbale della riunione stessa. 
Quando il problema viene segnalato, questo viene poi assegnato a uno o più membri del gruppo, i quali dovranno risolverlo e documentare la soluzione adottata.
Una volta risolto il problema, il responsabile di progetto deve controllare che la soluzione adottata sia corretta e che il problema non si ripresenti.

\subsubsection{Hotfix e correzioni minime}
Le correzioni minime sono interventi che non richiedono un'analisi approfondita del problema e possono essere risolti in modo rapido ed efficace. Queste modifiche possono essere apportate direttamente dai membri del team che individuano l'errore, senza necessità di assegnazione da parte del responsabile di progetto.  
Rientrano in questa categoria correzioni di errori di battitura, problemi di formattazione, imprecisioni sintattiche e altre anomalie di lieve entità. L'obiettivo delle correzioni minime è garantire un flusso di lavoro fluido e migliorare la qualità del progetto senza introdurre ritardi significativi. 

\subsubsection{Strumenti}
Qui vengono elencati gli strumenti usati, che vengono poi approfonditi nella sezione \hyperref[sec:strumenti]{Strumenti}.

\begin{itemize}
    \item GitHub;
    \item Zoom;
    \item Discord;
    \item Telegram;
    \item Meet;
\end{itemize}
